% from main.tex
%%%%%%%%%%  scan idx 265  |  folio 252  %%%%%%%%%%
supposer que \(X=Y[t_1,\dots,t_r]\), avec \(Y\) affine, \(f\) étant la projection canonique. \(Y\) est alors limite projective de schémas affines noethériens \(Y_\alpha\). Notons \(f_\alpha:X_\alpha=Y_\alpha[t_1,\dots,t_r]\to Y_\alpha\) la projection canonique. Il existe (EGA IV 8.5.2 et 11.2.6) un indice \(\alpha\) et un \(\mathcal O_{X_\alpha}\)-Module \(F_\alpha\) de présentation finie et \(f_\alpha\)-plat tel que \(g_\alpha^*(F_\alpha)\xrightarrow{\sim}F\), où \(g_\alpha:X\to X_\alpha\) est le morphisme canonique. Comme \(X_\alpha\) est noethérien, \(F_\alpha\) est cohérent, donc pseudo-cohérent rel. à \(f_\alpha\). Comme d'autre part \(F_\alpha\) est \(f_\alpha\)-plat et que \(f_\alpha\) est plat, on a \(g_\alpha^*(F)\xrightarrow{\sim}F_\alpha\otimes^{\mathbf L}_{Y_\alpha}\mathcal O_Y\), et il résulte de (1.10 a)) que \(F\) est pseudo-cohérent rel. à \(Y\), donc pseudo-cohérent, cqfd.

\textbf{Corollaire 4.3.1.} Soient \(f:X\to Y\) un morphisme localement de type fini, et \(E\) un \(\mathcal O_X\)-Module. Pour qu'on ait \(\parfamp_f(E)\subset[0,0]\), il faut et il suffit que \(E\) soit, rel. à \(f\), plat et de présentation finie (1.2.1 c)). En particulier, pour que \(f\) soit de dimension parfaite nulle (4.1), il faut et il suffit que \(f\) soit plat et localement de présentation finie.

\emph{Preuve.} Appliquer (4.3 (ii) \(\Longrightarrow\) (i)) avec \(a=b=0\).

\textbf{Corollaire 4.3.2.} Soient \(f:X\to Y\) un morphisme localement de type fini, et \(E\in\Ob D(X)\) un complexe de \(\mathcal O_X\)-Modules plats et de présentation finie rel. à \(f\) (1.2.1 c)), et \([a,b]\) un intervalle de \(\mathbb Z\). Si \(E^i=0\) pour \(i\notin[a,b]\), on a \(\parfamp_f(E)\subset[a,b]\). Si \(E\) est à degrés bornés supérieurement, \(E\) est pseudo-cohérent rel. à \(f\).

\emph{Preuve.} Conséquence de (4.3.1) et (I 4.15 et 2.11).

\textbf{Proposition 4.4.} Soit un triangle commutatif de schémas
\[
\begin{tikzcd}
X \arrow[r,hook,"i"] \arrow[d,"f"'] & Y \arrow[dl,"g"]\\
S &
\end{tikzcd}
\]
